\section{Evaluation}\label{sec:setup}
\subsection{Scenarios and kinematic tier}
Missions are created by a seeded generator (Supplementary Table~\ref{S-tab:params}). The map is a square with a side
of 10\,km and the depot at its centre. Half of the tasks are clustered around three hotspots, and half of them are
known at the start while the others arrive in batches. Time windows have a slack of 180 to 600\,s and services last 30
to 90\,s. Task rewards are 1, 2, 3 or 5 with probabilities 0.4, 0.3, 0.2 and 0.1, so the mean reward is 2.1. Three
fleet scales with 5, 10 and 20 UAVs and 25, 50 and 140 tasks are combined with three event profiles, namely six arrival
batches, three batches with two UAV failures, and three batches with four tightened deadlines.

The kinematic tier, called tier B, executes each published plan in continuous mission time. Legs follow a trapezoidal
speed profile with an acceleration limit, the airspeed of each leg deviates from the nominal speed by a multiplicative
error, a constant wind of random direction changes the ground speed, and service durations are scaled by a random
factor. Table~\ref{tab:levels} lists the three disturbance levels and the two latency levels. The planner never sees
these disturbances and receives the state of each UAV as telemetry. All random draws use common random
numbers~\cite{law2015simulation} keyed to the mission seed, the UAV and the task, so two planners that send the same
UAV to the same task meet the same disturbance.

\begin{table*}[t]
\caption{Disturbance and latency levels of the kinematic tier.}\label{tab:levels}
\centering\small
\setlength{\tabcolsep}{4pt}
\begin{tabular}{lcccc@{\quad}lccc}
\toprule
Execution & Acceleration (m/s$^2$) & Speed c.v. & Wind (m/s) & Service c.v. & Latency & Median (s) & Log-s.d. & Loss / timeout \\
\midrule
Low & 3.0 & 0.03 & 2 & 0.05 & Low & 0.2 & 0.5 & -- \\
Mid & 2.0 & 0.07 & 5 & 0.15 & High & 2.0 & 0.8 & 5\% / 5\,s \\
High & 1.5 & 0.12 & 8 & 0.25 & & & & \\
\bottomrule
\end{tabular}
\end{table*}

\subsection{Software-in-the-loop tier}
The software-in-the-loop tier, called tier C, flies missions with three X500 quadrotors, each controlled by its own
PX4~v1.17 autopilot~\cite{meier2015px4}, in Gazebo Harmonic~\cite{koenig2004gazebo} on ROS~2
Jazzy~\cite{macenski2022ros2}. Fig.~\ref{fig:tierc} shows how a mission is prepared and executed. A tier C mission is a
tier B scenario with three UAVs and fifteen tasks, rescaled by 0.1 in length and 0.25 in time. Speeds, energy, the
distance price and the latency factor are scaled accordingly, so every dimensionless ratio of the decision problem is
preserved. The mission manager runs the scenario timeline, builds the planning snapshot, calls the planner and releases
the plan. The executor streams offboard setpoints to every autopilot at 10\,Hz through the uXRCE-DDS bridge, flies
each leg towards a target point that moves along the straight leg, and holds the UAV at the task for the service
duration. The commanded ground speed follows the speed error and the scaled wind of tier B with the same random numbers,
while acceleration, attitude dynamics, position tracking and communication come from the simulated vehicles. Gazebo
runs in lockstep with PX4 at twice real time, and mission time is simulation time.

\begin{figure*}[t]
\centering
\fig{fig_tierc}{\linewidth}
\caption{Software-in-the-loop tier. Top: rescaling of a kinematic mission. Bottom: software
architecture.}\label{fig:tierc}
\end{figure*}

\subsection{Compared methods}\label{sec:methods}
Table~\ref{tab:methods} lists the compared methods. All of them share the repair, the ALNS operators and parameters,
the initial plan and the random seed, and every call has a hard budget of 250\,ms. They differ in the planning model,
the fallback, the rule that ends the search and the trigger. Unless stated otherwise, every method searches for a fixed
40\,ms, the mean computation of the stopping rule on development data, and replans periodically.

\begin{table*}[t]
\caption{Compared methods.}\label{tab:methods}
\centering\footnotesize
\setlength{\tabcolsep}{4pt}
\begin{tabular}{@{}>{\raggedright\arraybackslash}p{3.8cm}>{\raggedright\arraybackslash}p{4.6cm}%
>{\raggedright\arraybackslash}p{2.8cm}>{\raggedright\arraybackslash}p{4.1cm}c@{}}
\toprule
Method & Planning deadline & Fallback & Search & Studies \\
\midrule
Nominal model~\cite{ropke2006adaptive} & deadline & none & repair only, or ALNS with a fixed, scaled or stagnation budget & 1--7 \\
Conventional planner & deadline, event trigger only & none & ALNS, 40\,ms & 2--5 \\
AnytimeSwarm, AnytimeSwarm-R & deadline & none & stopping rule & 1, 2 \\
Performance profile~\cite{boddy1994deliberation,zilberstein1996using} & deadline & none & stopping rule & 2 \\
Online prediction~\cite{svegliato2018meta} & deadline & none & stopping rule & 2 \\
Constant buffer~\cite{jabali2015selfimposed} & deadline minus 60\,s & none & ALNS, 40\,ms & 2--5, 7 \\
Proportional buffer & deadline minus a tuned share of the lead time & none or nominal & ALNS, 40\,ms & 4, 5, 7 \\
Moment propagation~\cite{ehmke2015ensuring} & mean plus $z$ standard deviations along the route & none, nominal or cost-aware & ALNS, 40\,ms & 4--7 \\
CBBA~\cite{choi2009consensus,johnson2012nonsubmodular} & deadline or \mname{} deadlines & none or cost-aware & consensus & 6 \\
\mname{} (proposed) & effective deadline, Eq.~\eqref{eq:cc} & none, nominal or cost-aware & ALNS, 40\,ms or stopping rule & 2--7 \\
\bottomrule
\end{tabular}
\end{table*}

The nominal model plans against the true deadlines. Besides 40\,ms, it is run with regret-2 repair only, with 250\,ms,
with 1\,ms per open task and with a rule that stops after 20\% of the budget without improvement. The two stopping rules from the metareasoning literature use the same delay cost as
\mname{} and differ only in how they predict the gain of further search. The constant buffer, the simplest form of a time buffer in the
schedule~\cite{jabali2015selfimposed}, subtracts 60\,s from every deadline, the best of 30, 60 and 120\,s on
development missions. The proportional buffer subtracts a tuned share of the
lead time, so it keeps the shape of the \mname{} margin but replaces the calibration by one slope. Moment propagation
follows Ehmke et al.~\cite{ehmke2015ensuring} with leg parameters fitted on development flights. CBBA is
re-implemented with the bid warping of Johnson et al.~\cite{johnson2012nonsubmodular}, which time windows require for
convergence, on a fully connected network, and it pays for its planning time like every other method. Every model has
one parameter per disturbance level, selected on development missions. For \mname{}, the safety factor is 0.5 under
mid and 0 under high disturbances without fallback, and 1 and 0.5 with fallback. Supplementary
Section~\ref{S-sec:Sparams} gives the remaining parameters.

When a late arrival has a cost, a run is scored by
\begin{equation}\label{eq:score}
u_\lambda=U-\lambda\,\frac{n_{\mathrm{late}}}{n},\qquad u^w_\lambda=U-\lambda\,\frac{R_{\mathrm{late}}}{R},
\end{equation}
where $U$ is the utility, $n_{\mathrm{late}}$ of the $n$ tasks are served late, and they carry the reward
$R_{\mathrm{late}}$ of the total reward $R$. The score $u_\lambda$ charges every late arrival $\lambda$ mean task
rewards, and $u^w_\lambda$ charges $\lambda$ times the reward of the task. With $\lambda=0$, both equal the utility. The
competitors of the cost-aware fallback, namely the proportional buffer and \mname{} without fallback, are tuned on
development missions for every late cost and score, whereas the cost-aware fallback keeps the safety factor of the
utility.

\subsection{Studies and statistical protocol}\label{sec:protocol}
The evaluation consists of seven studies. Each study has its own held-out seeds, a written protocol and frozen code,
and each study after the first was designed after the held-out results of the earlier ones. Study~1 made only the
stopping rule execution-aware. Study~2 compared \mname{} without fallback with the baselines, study~3 added the nominal
fallback, and study~4 compared planning models on the same search. Study~5 introduced the cost-aware fallback, study~6
compared it with CBBA and applied it to moment propagation, and study~7 tested the late-cost model of the rule (7a) and
execution errors that the calibrated model does not describe (7b). Studies 5 and 6 ran the rule with a cost
proportional to the reward but scored it with $u_\lambda$, which charges a fixed cost per late arrival. This mismatch
between the cost model of the rule and the score was noticed after both studies had been frozen and run, so it could
not be corrected without changing their protocols. Study~7a was designed to test it on new seeds. It runs both cost
models, each scored with its own score. Study~7b draws speed and service errors from a Student-$t$ law, adds rare long delays, adds a persistent speed
bias per UAV, or replaces the latency by a heavy-tailed Pareto delay. Tier C flights belong to studies 2 and 4, the
latter called study~4C.

Development used seeds 0--9 in tier B and 0--2 in tier C. Before each study, the protocol with its hypotheses, the
configurations, the models and the code were frozen with SHA-256 hashes and committed to the version history, and the
held-out runs checked these hashes before starting. Studies 2 to 7 share one tier B matrix of two disturbance levels
(mid and high), high latency, three latency factors (10, 50 and 200) and two time-window factors (original and halved).
Each study uses 30 new held-out seeds. Every seed generates nine missions, one for each of the three fleet scales and
three event profiles, so a study runs 270 missions in twelve configurations. Study~4C flies \val{GN} missions under high
disturbances with halved windows, where tier B predicted the largest effect, with a size computed beforehand and one
interim analysis with O'Brien--Fleming bounds. A tier C run whose real-time factor leaves the range 1.6 to 2.2 is
repeated once.

The difference between two methods is averaged per mission over the configurations and tested with a paired
bootstrap~\cite{efron1993bootstrap} of 10\,000 resamples over missions. The primary hypotheses of each study form one
family, fixed in its protocol, and are corrected with the Holm procedure~\cite{holm1979simple}. Non-inferiority uses a margin of 0.5 percentage points (pp) and is followed by a
superiority test in a fixed sequence. All other comparisons are descriptive. Because the missions of one seed share
their execution disturbances, every tier B primary test was repeated with a bootstrap over seeds (Supplementary
Table~\ref{S-tab:Scluster}), which changed only the outcome of H9.
